\section{Introduction}
\label{sec:introduction}

Polynomial identity testing asks whether an arithmetic expression computes
the zero polynomial. We study this question for noncommuting variables when
the expression is accessible through matrix evaluations. The aim is to choose
the evaluation matrices before seeing the expression or any of its coefficients.

Let $F\langle x_1,\ldots,x_n\rangle$ be the free associative polynomial algebra
over a field $F$. A \emph{noncommutative formula} is a rooted tree whose leaves
are variables or scalars in $F$, and whose internal gates have two inputs and
are labelled $+$ or $\times$. The two inputs to each multiplication gate are
ordered. Its \emph{size} counts all gates, including the leaves. Write
$\mathsf{Form}_{n,s}(F)$ for the polynomials computed by formulas of size at
most $s$ in $n$ variables.

For a tuple $\mathbf T=(T_1,\ldots,T_n)\in\Mat_d(F)^n$, the value
$f(\mathbf T)$ is the usual ordered matrix evaluation; a scalar $a\in F$
acts as $aI_d$. A \emph{hitting point} for a class of polynomials is one tuple
on which every nonzero member of the class has nonzero matrix value.
More generally, a hitting set allows several tuples and requires a nonzero
value on at least one of them.

\begin{theorem}\label{thm:main}
There is a deterministic algorithm which, given integers $n,s\ge1$,
constructs a tuple
\[
 \mathbf T_{n,s}\in\Mat_d(\Q)^n,
 \qquad d\le 2ns^2,
\]
in time polynomial in $n,s$, measured in bit operations. The numerators and
denominators of its entries have bit length
$O(\log(ns+1)+ns^2\log(n+1))$.
For every field $F$ of characteristic zero and every
$0\ne f\in\mathsf{Form}_{n,s}(F)$,
\[
 f(\mathbf T_{n,s})\ne0,
\]
where the tuple is viewed over $F$ through the canonical embedding
$\Q\hookrightarrow F$. The construction uses $d=2ns^2$ when $ns\ge2$
and $d=1$ when $n=s=1$.
\end{theorem}

Thus a single exact matrix-evaluation query decides whether a polynomial
promised to lie in $\mathsf{Form}_{n,s}(F)$ is zero. The bit-time statement
concerns construction of the rational tuple. Evaluation over an arbitrary
field is understood as an exact oracle operation; no encoding or bit-cost
assumption on arbitrary elements of $F$ is needed. In particular, the singleton
set $\{\mathbf T_{n,s}\}$ gives a positive solution to the polynomial-size
rational-matrix hitting-set problem for characteristic-zero formulas.

\subsection{Context and scope}

An \emph{acyclic algebraic branching program} is a directed acyclic graph
with a designated source and sink, and with affine-linear edge labels
over the coefficient field. Its output is the sum of the ordered products
of edge labels along source-to-sink paths; here its size is the number
of vertices. Formulas admit such path representations of linear size.

Linear algebra has long played a central role in noncommutative complexity.
Nisan characterized layered homogeneous algebraic branching-program complexity by
ranks of coefficient matrices \cite[Theorem~1]{Nisan1991}.
Raz and Shpilka gave deterministic polynomial-time identity tests for
noncommutative formulas and algebraic branching programs
\cite[Theorems~4 and~5]{RazShpilka2005}. These algorithms use the internal structure
of the input, which is the white-box model. In the black-box model,
Bogdanov and Wee gave a randomized matrix-evaluation identity test with
one-sided error and running time polynomial in a supplied degree bound
and the other input parameters
\cite[Theorem~4.2]{BogdanovWee2005}.
Forbes and Shpilka constructed deterministic
quasipolynomial-size matrix hitting sets for noncommutative algebraic branching
programs over sufficiently large fields
\cite[Corollary~4.7]{ForbesShpilka2013}, and hence for formulas.
Their field-size condition is automatic in characteristic zero.

Arvind, Chatterjee and Mukhopadhyay obtained quasipolynomial black-box
tests over $\Q$ for rational formulas of inversion height two
\cite[Theorem~1.1]{ArvindChatterjeeMukhopadhyay2022}, and subsequently for
unrestricted noncommutative rational formulas
\cite[Theorem~2 and Corollary~3]{ArvindChatterjeeMukhopadhyay2026}.
Such formulas may have inverse gates, and evaluation requires each inverse
operand to be invertible. A separate manuscript
\cite[Theorem~1.1]{OpenAIRationalFormulas2026} constructs polynomial-size
rational matrix hitting lists for formulas over $\Q$ with a nonempty
rational matrix evaluation domain. Whenever such a formula is nonzero
on its domain, the list contains a defined evaluation with invertible
value, without a bound on inverse nesting or rational constant heights.
Here the single tuple works over every characteristic-zero coefficient
field for division-free formulas, whose matrix evaluations are everywhere
defined.

Two recent results treat different size restrictions. Lakhani and Mukhopadhyay
construct one integer matrix tuple for sparse noncommutative polynomials
over every characteristic-zero field, with polynomial construction time
and dimension equal to the sparsity bound, without a degree restriction
\cite[Theorem~2.1 and Remark~2.6]{LakhaniMukhopadhyay2026}.
Formula size can be much smaller than sparsity:
the ordered product of $k$ copies of $x_1+x_2$ has $4k-1$ gates and $2^k$
monomials. Lakhani and Saxena obtain randomized black-box tests for homogeneous
noncommutative circuits of constant depth
\cite[Theorem~1]{LakhaniSaxena2026}. Theorem~\ref{thm:main} is deterministic
and imposes no depth or homogeneity restriction on a formula.

The essential quantitative assertion is that matrix dimension, rational entry
length, and construction bit time are all polynomial in formula size and
the number of variables. For polynomials, a finite matrix hitting set can
always be combined into one tuple by taking block-diagonal sums, but the
resulting dimension is the sum of the dimensions in the set.
The construction here gives the polynomial
bounds directly. The formula model is used to obtain a small acyclic path
representation. In fact, Proposition~\ref{prop:path-programs} applies the
same construction to the algebraic branching programs defined above,
with matrix dimension $(n-1)\binom m2+m$ for $m$ vertices.
Thus the path-program consequence also improves the preceding
quasipolynomial hitting-set bound in characteristic zero.
No bound for arbitrary shared circuits is asserted.

\subsection{Proof idea}

The obstacle is that a small formula may contain exponentially many words
after expansion. We retain these words through a faithful encoding, and
bound the necessary truncation using a small representation of the formula.
We encode words by iterated integrals in one variable. For a word $w=i_1\cdots i_k$ on
$\{1,\ldots,n\}$, write $x_w=x_{i_1}\cdots x_{i_k}$, with
$x_\epsilon=1$ for the empty word. Define
\[
 h_{\epsilon}(z)=1,
 \qquad
 h_{iw}(z)=\int_0^z\frac{h_w(t)}{t+i}\,dt,
\]
where $\epsilon$ is the empty word and integration is interpreted by formal
power series at zero. Section~\ref{sec:iterated-integrals} proves that the
$h_w$ are linearly independent over every characteristic-zero field. Over
$\C$, continuation around the distinct poles yields a triangular spanning
argument in a truncated free algebra. A nonzero rational coefficient minor
then transfers independence to every field in the theorem. These functions
belong to the classical theory of iterated integrals and hyperlogarithms
\cite{Chen1977}. Their independence is a special case of the stronger
differential-field theorem of Deneufch\^atel, Duchamp, Hoang Ngoc Minh and
Solomon \cite[Theorem~1]{DeneufchatelEtAl2011}; we include a direct proof
for the particular kernels used here. Its triangular monodromy argument
has a two-letter antecedent in Minh, Petitot and van der Hoeven
\cite[Section~5, Theorem~4]{MinhPetitotVanderHoeven1998}.

If $f_w$ denotes the coefficient of $x_w$ in $f$, put
$h_f=\sum_w f_wh_w$.
Section~\ref{sec:formula-systems} converts a size-$s$ formula into an acyclic
path representation with at most $2s$ states, using the standard formula-to-path
construction of Nisan \cite[Section~2, Lemma~1]{Nisan1991}, also presented
by Raz and Shpilka \cite[Section~2.1, Lemma~2, Step~1]{RazShpilka2005}.
It follows that $h_f=uGv$ for
constant vectors $u,v$ and an invertible formal matrix $G$ satisfying
\[
 G'=\left(\sum_{i=1}^n\frac{B_i}{z+i}\right)G.
\]
The matrices $B_i$ describe the formula and are used only in the proof.

The quantitative step is a uniform bound on vanishing at zero.
Moura's multiplicity estimate for rational differential systems
\cite{Moura2003}, also stated in \cite[Section~2, Theorem~1]{MouraHyperelliptic2003},
bounds the order of a nonzero scalar component $uGv$
of an $m$-dimensional system $G'=(P/q)G$, with $q(0)\ne0$,
$\deg q=n$ and $\deg P\le n-1$, by
$(n-1)\binom m2+m-1$. Section~\ref{sec:multiplicity} gives a self-contained
formal proof of this estimate over the characteristic-zero field in use.
It uses the formal Wronskian criterion
\cite{BostanDumas2010} and the minimum order of the maximal minors of
a matrix of derivative rows. This bound is independent of all coefficient
values and applies beyond the particular systems arising from formulas.

Finally, Section~\ref{sec:construction} realizes the integration operators
on the space of Taylor coefficients of degrees below $2ns^2$.
Their matrices have explicit rational entries. Linear independence makes
$h_f$ nonzero, and the order bound ensures that this finite truncation retains
a nonzero coefficient. That coefficient occurs in the first column of
$f(\mathbf T_{n,s})$.

For $ns\ge2$, write $\pi_d(a)$ for the column of the first $d=2ns^2$
Taylor coefficients of a series $a$, and put $e_0=\pi_d(1)$.
Figure~\ref{fig:proof-route} summarizes the argument for a nonzero
formula polynomial $f$, separating the generator from the
formula-dependent objects used to prove its guarantee.
\begin{figure}[ht]
\centering
\begin{tikzpicture}[
  >=Latex,
  every node/.style={font=\small,align=center},
  every path/.style={line width=.45pt}
]
\node (bounds) at (0,1.55) {$n,s$};
\node (tuple) at (10.8,1.55) {$\mathbf T_{n,s}\in\Mat_d(\Q)^n$};
\draw[->] (bounds) -- node[above] {explicit rational entries; $d\le2ns^2$} (tuple);
\node (formula) at (0,0) {$f$};
\node (paths) at (3.2,0) {$(u,B_1,\ldots,B_n,v)$};
\node (series) at (7.05,0) {$h_f=uGv\ne0$};
\node (column) at (10.8,0) {$f(\mathbf T_{n,s})e_0\ne0$};
\draw[->] (formula) -- node[above] {paths} (paths);
\draw[->] (paths) -- node[above] {word integrals} (series);
\draw[->] (series) -- node[above] {$\pi_d$} (column);
\draw[->] (tuple) -- (column);
\node[font=\footnotesize] at (5.4,-.65)
  {The lower row is used only to prove the hitting property.};
\end{tikzpicture}

\caption{The construction for $ns\ge2$. The generator uses only the size bounds.
The correctness argument
associates a differential system to each possible input formula. Independence
makes its scalar output nonzero, and the multiplicity bound ensures that
the retained Taylor coefficients include a nonzero entry.}
\label{fig:proof-route}
\end{figure}
